\documentclass[11pt]{article}
\usepackage[margin=1.1in]{geometry}
\usepackage[cmex10]{amsmath}
\usepackage{amssymb,amsthm}
\usepackage[T1]{fontenc}
\usepackage[hidelinks,hypertexnames=false]{hyperref}
\allowdisplaybreaks[2]
\renewcommand{\labelitemi}{$\bullet$}
\setlength{\emergencystretch}{2em}
\title{Eventual negative zeros for balanced arrays}
\author{ProgramComputer}
\date{}
\newtheorem{theorem}{Theorem}
\newtheorem{corollary}[theorem]{Corollary}
\DeclareFontShape{OMX}{cmex}{m}{n}{<-> cmex10}{}
\begin{document}
\maketitle

\begin{abstract}
We study the polynomial counting labeled binary arrays with equal column
sums and a bound on every row sum. The theorem gives a row
threshold uniform in the dimension and cap, above which all zeros are
simple and strictly negative. The supporting estimates concern angular
signs and the probability of exact column balance for prescribed sparse
row sizes.
\end{abstract}

\section{Statement}
The \href{https://www.symmetricfunctions.com/realRootedOpenProblems.htm\#conjBalancedArrays}{balanced-array conjecture of Per Alexandersson}
asks for real-rootedness throughout the parameter range. Write
$P_n^{(r,d)}(q)=\sum_k a_kq^k$, where $a_k$ counts labeled binary
$n\times d$ arrays with every column sum $k$ and every row sum at most $r$.

\begin{theorem}
There is an absolute $N$ such that for all integers $n\ge N$, $d\ge2$,
and $1\le r\le d$, $P_n^{(r,d)}$ has exactly $\lfloor nr/d\rfloor$
simple, strictly negative zeros.
\end{theorem}
For $d=1$ the polynomial is $(1+q)^n$. The theorem leaves a
finite range of possible exceptions to the original conjecture, with an ineffective row cutoff, by
\href{https://github.com/ProgramComputer/open-math-research/blob/78e9ef2c0c0b7a5daed126efbe12a46d030d8200/topics/balanced-arrays/orientation-normalization.md}{coefficient separation}
for $d\ge n$.

\section{Reduction}

Suppose the theorem fails for arbitrarily large row counts and choose a sequence of failures with $n\to\infty$. We use two previously established bounds. For every fixed $\alpha>0$, once $n$ is sufficiently large in terms of $\alpha$, all caps are covered at $n\ge\alpha d^2$. Also an absolute $B$ covers every cap at $d\ge B n/\log n$ for all large $n$. Full rank has simple negative roots in every dimension $d\ge2$. Passing to subsequences as necessary, the remaining parameters satisfy

\begin{equation*}
d\to\infty,\qquad 1\le r<d,\qquad
b=d^2/n\to\infty,\qquad b/d=d/n\to0 .
\end{equation*}

In particular $n/d\to\infty$, $n<d^2$ eventually, and $\log n=o(d)$. We will exclude every such sequence using only these relations. \href{https://github.com/ProgramComputer/open-math-research/tree/605cc467a5be93a8408519da331e46f1008df96e/topics/balanced-arrays/critical-aspect}{Fixed quadratic aspect} \enspace \href{https://github.com/ProgramComputer/open-math-research/tree/4fb624ba937a8eecc1361a6066185cd81ab7b172/topics/balanced-arrays/dense-regime}{Large-dimension bound} \enspace \href{https://github.com/ProgramComputer/open-math-research/blob/4fb624ba937a8eecc1361a6066185cd81ab7b172/topics/balanced-arrays/finite-linear-reduction.md}{Full rank}

Every coefficient in the displayed support is positive: assign $dk$ slots in row groups of size at most $r$ and fill them consecutively with the repeated column cycle $1,\ldots,d$. Each group then has distinct columns. Write $T=nr/d$, $t_k=\sqrt{a_{k-1}/a_{k+1}}$ and $\Gamma_k=a_k^2/(a_{k-1}a_{k+1})$ for internal $k$.

\section{Counting at bounded column quota divided by dimension}

Pass to a subsequence along which $r/b$ converges in $[0,\infty]$. If $r/b$ is bounded, $D/d$ is bounded, and every internal geometric radius $t_k$ gives $(-1)^kP(-t_k)>0$ eventually. This is the bounded-rank-ratio lemma, whose coefficient comparison now holds throughout $b=o(d)$. It conditions on real row sizes and uses the fixed-profile balance estimate

\begin{equation*}
\mathbb P(\mathrm{balance}\mid s_1,\ldots,s_n)
=\sqrt d\,(2\pi\Lambda)^{-(d-1)/2}
\exp\left(-\frac1{12\beta}\right)(1+o(1)),\qquad
\Lambda=\sum_i \frac{s_i(d-s_i)}{d(d-1)} .
\end{equation*}

Its hypotheses are $\sum_i s_i=dk$, $k/d\to\beta\in(0,\infty)$ and $\max_i s_i=o(d)$. The written proof of this estimate includes its global localization and the Gaussian cubic and quartic corrections. To average over rows at fixed total, the dependence of the determinant on their empirical variance is retained. The tilted row law has a Poisson exponential-moment bound, which controls that variance even when the largest permissible degree grows almost linearly in $d$. These estimates are also used at the two ends in the remaining cases. \hyperref[sec:sparse-fixed-profile-edgeworth]{Balance estimate} \enspace \hyperref[sec:sparse-fixed-profile-global-localization]{Localization for fixed sizes} \enspace \hyperref[sec:edgeworth-bounded-rank-ratio]{Conditional averaging and bounded-rank signs}

The orientation inequality gives strict log-concavity of $a_k$ at internal quotas and $\Gamma_k\ge(1+1/k)^d(1+1/(n-k))^d$. Thus the geometric radii are strictly ordered. If $D\le1$ the conclusion is immediate. Otherwise the bounded-rank signs and the endpoint signs give $D$ distinct roots, excluding this case. \href{https://github.com/ProgramComputer/open-math-research/blob/4fb624ba937a8eecc1361a6066185cd81ab7b172/topics/balanced-arrays/orientation-normalization.md}{Orientation inequality}

\section{The phase radii}

We may now assume $r/b\to\infty$ and thus $D/d\to\infty$. Given $w>0$ let $S$ have the positive size law with weights $\binom ds w^s$, $0\le s\le r$, and let $H(w)$ be their sum. Define $\mu=\mathbb ES$, $h=r-\mu$, $y=n\mu/d$, $G(s)=s(d-s)/(d(d-1))$, $a=\mathbb E G(S)$, $\phi=\pi/d$ and $\nu=(d-1)/2$. The occupied-cell phase is

\begin{equation*}
\Psi(w)=n\arg H(we^{i\phi})-
\nu\arg\frac{\mathbb E[G(S)e^{i\phi S}]}{\mathbb E e^{i\phi S}} .
\end{equation*}

The arguments are obtained by continuity from zero angle since both sums have positive weights and angles in a sector of width $r\phi<\pi$. The limits at radii zero and infinity are $-\nu\phi$ and $T\pi$. For each $1\le k<D$ the first value $v_k>0$ where $\Psi(v_k)=k\pi$ exists and these first hitting points strictly increase. The tilt inequality

\begin{equation*}
|\Psi(w)/\pi-y|\le C d+C\min(y,T-y)/\sqrt d
\end{equation*}

holds throughout, with an absolute $C$. Therefore at $w=v_k$, $k/d\to\infty$ implies $\mu/b=y/d\to\infty$, and $(D-k)/d\to\infty$ implies $h/b=(T-y)/d\to\infty$. The variance inequality $\mathrm{Var}(S)\le\mu(1-\mu/d)$ gives

\begin{equation*}
\frac{na}{d}\ge\frac{\mu(d-\mu)}{bd}
\ge\frac{\min(\mu,d-\mu)}{2b}.
\end{equation*}

\href{https://github.com/ProgramComputer/open-math-research/blob/4fb624ba937a8eecc1361a6066185cd81ab7b172/topics/balanced-arrays/quadratic/proof.md}{Phase and variance bounds}

At any internal first hitting point with $\mu\ge1/10$ and $na/d\to\infty$ the angular lemma gives $(-1)^kP(-v_k^d)>0$. The negatively evaluated polynomial is the torus integral over coordinate product $-1$ of the $n$th power of the row generating function. Uniformly controlling its angular logarithm uses a zero-free disk for each fixed-size row, a consequence of stability of the elementary symmetric polynomial. Averaging the fixed-size logarithmic estimates after subtracting their common quadratic and cubic parts controls the fluctuations due to the random row size. This removes a previous restriction on the relative growth of $n$ and $d$. Near the last internal level, more than two angular regions can contribute. Their significant odd phases have the same sign. \hyperref[sec:analytic-angular-signs]{Angular sign theorem}

\section{The overlap and the full degree}

At $k/d\to0$, the orientation lower bound gives $\Gamma_k\to\infty$. The sign at $t_k$ is $(-1)^k$, since after accounting for the two neighboring terms the signed central sum is at least $1-2/\sqrt{\Gamma_k}$ by alternation of the decreasing tails. At $k/d\to\beta\in(0,\infty)$ the lower-quota overlap lemma gives $\log\Gamma_{k+j}\to1/\beta$ at every fixed offset and $d\log v_k=\log t_k+o(1)$. The signed sum at both radii tends after normalization by its central term to $\sum_{j\in\mathbb Z}(-1)^j e^{-j^2/(2\beta)}>0$ by Gaussian Poisson summation. Log-concavity controls the tails from the ratios at any fixed large offset. \hyperref[sec:edgeworth-finite-lower-overlap]{Lower-quota overlap}

Pass to a further subsequence on which the corank ratio $(d-r)/b$ tends to infinity or to $\gamma\in[0,\infty)$. In the first case use the radii $t_k$ for $k\le d$ and $v_k^d$ for $k>d$. For $k/d\to\infty$, we have $\mu/b\to\infty$ and $(d-\mu)/b\ge(d-r)/b\to\infty$. Therefore $na/d\to\infty$ at the first hitting point and its sign follows from the angular lemma. This covers the entire upper end regardless of the slack.

In the second case use $t_k$ also when $D-k\le d$. The two geometric end ranges are disjoint for large $d$. Whenever both $k/d$ and $(D-k)/d$ tend to infinity, the angular lemma again applies since $\mu/b\to\infty$ and $(d-\mu)/b\ge h/b\to\infty$. For $(D-k)/d$ bounded, apply the complementary quota lemma to the holes: the hole column sum is $n-k$ and each row must have at least $d-r$ holes. This gives the strict sign at the geometric radius even when $(D-k)/d\to0$, including an exactly saturated last coefficient as neighbor. When $(D-k)/d\to\beta\in(0,\infty)$ it gives $\log\Gamma_{k+j}\to1/(\gamma+\beta)$ and $d\log v_k=\log t_k+o(1)$, so either radius has the required sign. A switch moving a single hole bounds the upper tail of row sizes under the actual balanced count. This permits use of the fixed-profile estimate for typical sizes even though the hole rows have no prescribed small maximum. \hyperref[sec:sublinear-b-generic-upper-transfer]{Complementary quota lemma}

All radii increase strictly inside their respective geometric and angular ranges. At the lower join the angular substitution is at $d+1$. When there is an upper geometric block, the substitution is at $D-d-1$, immediately before the first geometric point at $D-d$. At either join the adjacent geometric log radii differ by $(\log\Gamma_k+\log\Gamma_{k+1})/2$. The limit is $1$ at $k$ near $d$, and $1/(\gamma+1)>0$ near $D-d$ when this second join is used. Replacing a geometric radius by its first hitting radius there changes its log by $o(1)$. Thus the full family is ordered. If a sign failed infinitely often, a subsequence would have both $k/d$ and $(D-k)/d$ convergent in $[0,\infty]$. Their sum tends to infinity. The cases above contradict every possible failure. With the signs $1$ at zero and $(-1)^D$ at positive infinity, the $D-1$ internal alternating signs give $D$ distinct negative zeros. This contradiction proves that a single integer $N$ works for every dimension and cap.

For each of the finitely many $n<N$, all $d\ge n$, $d\ge2$ are already covered by coefficient separation and $d=1$ by the explicit formula. Hence only finitely many triples can fail the original conjecture. This proof gives no numerical bound on their row counts. \href{https://github.com/ProgramComputer/open-math-research/blob/4fb624ba937a8eecc1361a6066185cd81ab7b172/topics/balanced-arrays/finite-linear-reduction.md}{Coefficient separation in large dimensions}


\appendix
\clearpage
\section{Angular signs from a zero-free transverse expansion}\label{sec:analytic-angular-signs}

Let $d\to\infty$, $1\le r<d$, and use the capped-binomial row tilt $S$ with masses proportional to $\binom ds v^s$ for $v>0$. Write $\mu=\mathbb ES$, $h=r-\mu$, $V=\mathrm{Var}(S)$, $a_s=s(d-s)/[d(d-1)]$ and $a=\mathbb E a_S$. Assume $\mu\ge1/10$ and $na/d\to\infty$ (in particular $n/d\to\infty$). Set $\phi(t)=\mathbb E e^{it S}$, $a(t)=\mathbb E[a_S e^{itS}]/\phi(t)$, $t_*=\pi/d$, $\nu=(d-1)/2$, and $D=\lfloor nr/d\rfloor$. At every internal node

\begin{equation*}
n\arg\phi(t_*)-\nu\arg a(t_*)=k\pi,\qquad 1\le k\le D-1,
\end{equation*}

the original balanced-array polynomial satisfies $(-1)^kP_n^{(r,d)}(-v^d)>0$ eventually along every such parameter sequence. The arguments and first hitting points are those in \href{https://github.com/ProgramComputer/open-math-research/blob/78e9ef2c0c0b7a5daed126efbe12a46d030d8200/topics/balanced-arrays/fifth-thirds/subquadratic-angular-nodes.md}{the earlier phase-combination proof}. In particular this result includes first hitting points satisfying the stated mass condition. It allows arbitrary $nV/d^2$ and imposes no separate polynomial lower bound on $n$.

\subsection{Fixed-size analytic expansion}

Let $y\in\mathbf1^\perp$, $M=\|y\|_\infty$, $Q_j=\sum_i y_i^j$, $L_s=\sum_{i\in I}y_i$ for a uniform $s$-subset, and $K_s(z)=\log\mathbb E e^{zL_s}$ with its branch zero at $z=0$. The case $M=0$ is trivial. For $|z|M<c_0$ with absolute $c_0<\pi/2$, the arguments of $x_i=e^{zy_i}$ are in $(-\pi/2,\pi/2)$. Then $e_s(x)\ne0$: the polynomial $\prod_i(w+x_i)$ has all roots in the open left half-plane, so all roots of its derivatives stay there by Gauss--Lucas, and its derivative of order $d-s$ at zero is $(d-s)!e_s(x)$. Thus the specified analytic logarithm exists on that disk.

For real $u$ with $|u|M\le c_0$, Hoeffding's convex comparison for sampling without replacement, with $\exp(uL_s)$, bounds the mgf by that for $m=\min(s,d-s)$ iid population draws, reversing sign by complement when needed. Since $\sum_i y_i=0$, the one-draw mgf is at most $1+C u^2Q_2/d$, using the second-order exponential Taylor bound on this interval. As $m/d\asymp a_s$ for $1\le s<d$,

\begin{equation*}
\mathrm{Re}K_s(z)=\log|\mathbb E e^{zL_s}|
\le \log\mathbb E e^{(\mathrm{Re}z)L_s}
\le C a_s |z|^2Q_2.
\end{equation*}

The size zero contributes zero. Hoeffding's comparison and an elementary coupling proof are also recorded in \href{https://github.com/ProgramComputer/open-math-research/blob/78e9ef2c0c0b7a5daed126efbe12a46d030d8200/topics/balanced-arrays/fifth-thirds/subquadratic-interior-saddle.md}{the fifth-moment note}. \href{https://doi.org/10.1080/01621459.1963.10500830}{Hoeffding Theorem 4}.

For $R_0=c_0/(2M)$, Borel--Carathéodory, using the real-part bound and $K_s(0)=0$, bounds $|K_s|$ on $|z|\le R_0/2$ by $C a_s Q_2R_0^2$. Cauchy's coefficient bound there and summing a geometric series show that the Taylor tail at $z=i$ from degree five on is at most $C a_s Q_2M^3$, provided $M\le\eta$ for a sufficiently small absolute $\eta$. The quadratic and cubic coefficients at $i$ are $-a_sQ_2/2$ and $-i\gamma_s Q_3/6$, with $\gamma_s=a_s(d-2s)/(d-2)$. The fourth coefficient obeys $O(a_sQ_4+a_sQ_2^2/d)$: the $Q_2^2$ moment term divided by 3 is $(s)_2(d-s)_2/(d)_4=a_s^2+O(a_s/d)$, and the logarithm subtracts the $a_s^2$ term. Hence on a fixed small ball $\|y\|\le\eta$,

\begin{equation*}
K_s(i)=-a_sQ_2/2-i\gamma_sQ_3/6+E_s(y),\qquad
|E_s(y)|\le C a_s(Q_4+Q_2M^3),
\end{equation*}

using $Q_2^2/d\le Q_4$.

\subsection{Complex mixture and Gaussian integration}

As in \href{https://github.com/ProgramComputer/open-math-research/blob/4fb624ba937a8eecc1361a6066185cd81ab7b172/topics/balanced-arrays/quadratic/very-tall-quadratic-interior-saddle.md}{the quadratic interior note}, if $h\ge h_0>0$ keep near charts $|t|\le c/\sqrt{1+V}$; if $0<h\le h_0$ for sufficiently small fixed $h_0$, keep every chart, using \href{https://github.com/ProgramComputer/open-math-research/blob/4fb624ba937a8eecc1361a6066185cd81ab7b172/topics/balanced-arrays/quadratic/near-full-all-angular-signs.md}{the near-full estimates}. In both settings $|\phi(t)|$ is bounded below uniformly. Put $\mathbb E_t X=\mathbb E[Xe^{itS}]/\phi(t)$, $\gamma(t)=\mathbb E_t\gamma_S$ and write $R(t\mathbf1+y)=\mathbb E e^{itS+iL_S}$. To average the fixed-size expansion factor out $\exp(-a(t)Q_2/2-i\gamma(t)Q_3/6)$ and let $D_s=-(a_s-a(t))Q_2/2-i(\gamma_s-\gamma(t))Q_3/6$. Then $\mathbb E_t D_S=0$. The functions $a_s,\gamma_s$ are $O(1/d)$-Lipschitz, and the complex $t$ weights are dominated by a constant times their radial weights. Centering each function first at its value at $\mu$ gives
$\mathbb E|a_S-a(t)|^2+\mathbb E|\gamma_S-\gamma(t)|^2\le C V/d^2\le Ca/d$ because $a\ge V/d$.
On a small fixed ball, $D_s,E_s$ are uniformly small. Thus expanding $\mathbb E_t e^{D_S+E_S}$ about one and using the vanished linear $D_S$ term bounds its logarithm by $C(\mathbb E|D_S|^2+\mathbb E|E_S|)$. Since $|Q_3|\le M Q_2$, the $D_S$ bound is $CaQ_2^2/d\le CaQ_4$. We obtain

\begin{equation*}
\log(R(t\mathbf1+y)/\phi(t))
=-a(t)Q_2/2-i\gamma(t)Q_3/6+O\!\left(aQ_4+aQ_2M^3\right).
\end{equation*}

Let $A=n\Re a(t)\asymp na$ and put $B_t=(1+V)t^2$ in the near interior charts, and $B_t=h\min(t^2,1)$ in the near-full charts. The earlier coefficient bounds give $\Re a(t)/a=1+O(B_t)$ and $|\mathrm{Im}\gamma(t)|\le Ca\sqrt{B_t}$. For the last bound in the near-full regime one can use the stronger $Ca h\min(|t|,1)$ from its weighted deficit moments. The positive correction to the real Gaussian exponent is bounded by

\begin{equation*}
Cna\left(\sqrt{B_t}\sum_i |y_i|^3+Q_4+Q_2M^3\right).
\end{equation*}

The coordinate terms have bounded exponential expectation up to $\exp(O(d B_t+d/A))$ on the ball, also with any fixed multiplier. To see this, write $Y_i=Z_i-\bar Z$ for iid $N(0,1/A)$, impose $|\bar Z|\le\eta$ independently of $Y$, and use $|Z_i|\le2\eta$ and $\sum |Y_i|^j\le2^j\sum|Z_i|^j$. In the resulting one-dimensional factors the correction is absorbed by half the Gaussian quadratic on the fixed truncated interval. Integrating $e^x-1\le x e^x$ bounds each factor by $1+C(\sqrt{B_t/A}+1/A)$; the product bound follows by Young's inequality.

For the remaining $A Q_2M^3$ term, $U=A Q_2$ is $\chi^2_{d-1}$. On $U\le Ld$ for fixed large $L$, the exponent is at most $CLd M^3$. When $M\le x_0=\sqrt{K\log(d+1)/A}$, it is $O_L(d(\log(d+1)/A)^{3/2})=o(1)$ since $A/d\to\infty$. For $x_0<M\le\eta$, it is at most $A M^2/8$ for large $d$ by $A/d\to\infty$. The Gaussian union bound $\mathbb P(M>x)\le2d e^{-A x^2/2}$ and integration of that tail show that the expected exponential on this exceptional part is $O(d e^{-3A x_0^2/8})=o(1)$ for large fixed $K$. Finally on $U>Ld$ inside the chart, absorb $CU M^3$ by a fixed small weakening of the Gaussian precision. The $\chi^2$ exponential moment bounds this tail by $e^{-cd}$ for large $L$ and small $\eta$. This argument holds with any fixed multiplier, so Cauchy--Schwarz combines it with the coordinate bounds.

With $I_0=(2\pi n a)^{-\nu}/\sqrt d$, the absolute chart integral is thus at most $C I_0|\phi(t)|^n e^{Cd B_t}$. Its difference from $G_t=I_0\phi(t)^n(a/a(t))^\nu$ is at most the same bound times

\begin{equation*}
O\!\left(\sqrt{\frac d{na}}+\frac d{na}
+d\left(\frac{\log(d+1)}{na}\right)^{3/2}\right)=o(1).
\end{equation*}

Indeed under the real Gaussian $\|Q_3\|_2=O(\sqrt d A^{-3/2})$ by Wick's zero-sum cancellation, and $\|Q_4\|_2=O(d A^{-2})$. Also $\|Q_2 M^3\|_2\le \|Q_2\|_4\|M^3\|_4=O((d/A)(\log(d+1)/A)^{3/2})$ by the chi-square and maximum moment bounds. Multiplying by $na$ and using Cauchy--Schwarz for the difference of exponentials proves the error claim. Gaussian tails outside the fixed-radius chart are negligible since $A/d\to\infty$. The factor $1/\sqrt d$ is the normalized Haar covolume correction, as in the prior angular notes.

\subsection{Global localization}

Here is the additional modulus estimate needed when $n/d$ grows arbitrarily slowly. Put $z_i=e^{i\theta_i}$, $\bar z=d^{-1}\sum_i z_i$, $D_z=\sum_i |z_i-\bar z|^2=d(1-|\bar z|^2)$. At fixed size $s\ge1$, express its Fourier response as the Fourier coefficient for iid Bernoulli variables of parameter $p=s/d$, conditional on total $s$. The probability of that total is bounded below by $c/\sqrt{1+d p(1-p)}$ by the binomial mass and Stirling bounds. In the numerator bound the product of Bernoulli Fourier factors by

\begin{equation*}
\exp\!\big[-c p(1-p)d(1-|\bar z|)-c p(1-p)d|\bar z|(1-\cos u)\big],
\end{equation*}

after rotating the common integration variable $u$ by $\arg\bar z$. This follows from $|1-p+pe^{i\psi}|^2=1-p(1-p)|1-e^{i\psi}|^2$. If $|\bar z|\ge1/2$, integrating the $u$ factor costs $O(1/\sqrt{1+d p(1-p)})$, canceling the conditioning factor. If $|\bar z|<1/2$, then $D_z\asymp d$ and the prefactor $\sqrt{1+d p(1-p)}$ is absorbed for sufficiently large $a_sD_z$. Since $a_s\asymp p(1-p)$, for large $a_sD_z$ this proves $|f_s|\le e^{-c a_sD_z}$. For bounded $a_sD_z$, use the \href{https://github.com/ProgramComputer/open-math-research/blob/04564a728c21eab6ce511b98a48ccaac754648ff/topics/balanced-arrays/very-tall/very-tall-interior-saddle.md}{disjoint-pair bound} $|f_s|\le1-c\min(a_sD_z,1)$ and adjust constants. Hence the exponential bound holds for all sizes, including $s=0$.

To average, $a$ is comparable to $\mu(d-\mu)/d^2$ because $V\le\mu(1-\mu/d)$. The same variance inequality at every exponential tilt implies
$\mathbb P(S<\mu/2)\le e^{-c\mu}$ and $\mathbb P(d-S<(d-\mu)/2)\le e^{-c(d-\mu)}$. For example, for the first inequality the tilted mean $m(t)$ obeys $0\le m'(t)\le m(t)$, so $m(-t)\ge\mu e^{-t}$ for $t\ge0$. Integrating gives $\log\mathbb E e^{-tS}\le-\mu(1-e^{-t})$, and Chernoff with $t=\log2$ works. The same argument applies to $d-S$ by the variance bound at every tilt. On the complement of these events $a_S\ge c a$. Since $\min(\mu,d-\mu)\asymp da$ and $D_z\le d$, averaging the fixed-size bound gives $|R|\le3e^{-c'aD_z}$. For bounded $aD_z$, the earlier weighted disjoint-pair bound removes the prefactor: $\mathbb E a_S^2\le C a^2$ by Lipschitz continuity, $a\ge c/d$ and $a\ge V/d$, so $\mathbb P(a_S\ge a/2)\ge c$ by Cauchy--Schwarz. Consequently for all phases $|R|\le e^{-c'' aD_z}$.

Outside the fixed-radius charts $D_z\ge c\eta^2$. Retain half the penalty $e^{-cna\eta^2}$ and integrate the remainder by partitioning over closest anchors. Relative angle coordinates push forward normalized Haar measure to product Haar, by the onto torus homomorphism explained in the quadratic proof. Chord bounds give an integral at most $d(C/\sqrt{na})^{d-1}e^{-cna\eta^2}$. Relative to $I_0|\phi(t_*)|^n$, the logarithm is at most $O(d)-cna\eta^2+O(nV/d^2)$, tending to minus infinity since $V/d^2\le a/d$ and $na/d\to\infty$.

For $h\ge h_0$, far diagonal charts have a fixed modulus gap by the earlier diagonal comparison and $|R-\phi(t)|\le aQ_2/2$. Their aggregate measure is at most $d(C/\sqrt d)^{d-1}$: the Euclidean volume of each $(d-1)$-ball of fixed radius has this bound by Stirling, and Haar probability density is $1/[\sqrt d(2\pi)^{d-1}]$. Thus their relative contribution is at most $\exp[-cn+O(d(1+\log(na/d)))+O(nV/d^2)]$, which tends to zero since $n/d\to\infty$, $a\le1$, and $V/d^2=O(1/d)$.

\subsection{Diagonal reinforcement and sign}

The chart bounds and global tails now match those used in \href{https://github.com/ProgramComputer/open-math-research/blob/78e9ef2c0c0b7a5daed126efbe12a46d030d8200/topics/balanced-arrays/fifth-thirds/subquadratic-angular-nodes.md}{the all-angular sign proof} with a different $o(1)$ local error and no height assumption. For $h\ge h_0$, $V\ge c_{h_0}$, so $nV\gg d(1+V)$. When $nV/d^2\to\infty$, the nonprincipal near charts are suppressed by the diagonal gap $\exp[-c n V(t^2-t_*^2)]$, which absorbs $e^{Cd B_t}$. When $nV/d^2$ is bounded retain odd labels $|t|\le L/\sqrt{nV}$ for large fixed $L$. Subtracting $t/t_*$ times the principal phase cancels the linear terms in both arguments and leaves $O_L(\sqrt{(1+V)/(nV)}+d(1+V)/(nV))=o(1)$. Their real parts have the prescribed common sign because $t/t_*$ is odd. The discarded Gaussian lattice tail can be made arbitrarily small relative to the retained mass by taking $L$ large.

For $0<h\le h_0$, use the near-full gap with $h$ replacing $V$, valid at all labels, and $nh\gtrsim d$ at internal nodes by the near-full phase calculation. In the bounded $nh/d^2$ case the same retained phases at width $L/\sqrt{nh}$ have alignment error $O_L((nh)^{-1/2}+d/n)=o(1)$; when $nh/d^2\to\infty$ only the principal pair matters. The lattice tail estimate again uses only $n/d\to\infty$. These comparisons exhaust subsequences and give the sign asserted above.

\section{Global localization for sparse fixed-profile equal-column conditioning}\label{sec:sparse-fixed-profile-global-localization}

Let $d\to\infty$. For each row with prescribed size $s$ choose a uniform $s$-subset of the $d$ columns, independently across rows. Rows of size zero may be omitted. Suppose that the total size is $M=dk$, with $k/d\to\beta\in(0,\infty)$, and the maximal row size $s_{\max}=o(d)$. Write

\begin{equation*}
f_s(\vec z)=\frac{e_s(\vec z)}{\binom ds},\qquad
\Lambda=\sum_{\rm rows}\frac{s(d-s)}{d(d-1)}
       =k(1+o(1)),\qquad
G_d=\sqrt d\,(2\pi\Lambda)^{-(d-1)/2}.
\end{equation*}

Here each $z_i$ lies on the unit circle. For any fixed $\eta>0$, there is a fixed $L$ depending on $\beta$ such that the absolute contribution of the following two regions to the balanced-column integral is $o(G_d)$:

\begin{itemize}
\item $D_z=\sum_i|z_i-\overline z|^2>L$, where $\overline z=d^{-1}\sum_i z_i$.
\item $D_z\le L$ and some $z_i$ is at angular distance at least $\eta$ from $\overline z/|\overline z|$.
\end{itemize}

The estimates are exponential relative to $G_d$ for fixed $\eta$ and fixed $\beta$. Thus the remaining quotient chart has bounded centered sum of squared angles and arbitrarily small fixed maximum angle deviation. More generally one may take $\eta=\eta_d\downarrow0$ as long as $\eta_d^2\gg (s_{\max}/d)^2+(\log d)/d$, using the refinement at the end.

The character integral on the full torus with product Haar measure is

\begin{equation*}
\Pr(\text{column sums all }k)
=\int \left(\prod_i z_i^{-k}\right)\prod_{\rm rows} f_s(\vec z)\,d\vec z.
\end{equation*}

The integrand is invariant under simultaneous rotation because $M=dk$. In particular estimates on the full torus apply directly to its quotient integral. In the local zero-sum angular hyperplane, the projection of the chart with last angle fixed has Gram determinant $1/d$, so the normalized Gaussian measure of precision $\Lambda$ is $G_d$. The bounds below are absolute and therefore independent of the choice of character.

\subsection{A uniformly exponential bound for each row}

Write $D=D_z=d(1-|\overline z|^2)$. For every $1\le s\le d/2$,

\begin{equation}
|f_s(\vec z)|\le \exp(-c sD/d).
\tag{G1}
\end{equation}

First use the Cauchy coefficient circle of radius $r=s/(d-s)$ in $\prod_i(1+v z_i)$. Put $p=s/d$. Since $|1+r e^{i\theta}|/(1+r)\le \exp[-p(1-p)(1-\cos\theta)]$,

\begin{equation*}
|f_s(\vec z)|
\le
\frac{(1+r)^d}{r^s\binom ds}
\,e^{-p(1-p)d(1-|\overline z|)}
\int_{-\pi}^{\pi}e^{-p(1-p)d|\overline z|(1-\cos\theta)}
                    \,\frac{d\theta}{2\pi}.
\end{equation*}

The prefactor is the inverse mass of $\mathrm{Bin}(d,p)$ at its mean and is at most $C\sqrt s$ by Stirling. The integral is at most $C/\sqrt{1+s|\overline z|}$ by $1-\cos\theta\ge c\theta^2$ on this interval and $1-p\ge1/2$. If $|\overline z|\ge1/2$, the combined prefactor is bounded, and $1-|\overline z|\ge D/(2d)$. If $|\overline z|<1/2$, the exponent gives $e^{-cs}$, which absorbs $\sqrt s$ with a reduction of $c$. In both cases this proves $|f_s|\le C e^{-c sD/d}$.

To remove that constant when $sD/d$ is bounded, generate a uniform size-$s$ subset by taking $s$ disjoint random pairs from a uniform permutation and independently choosing one element of each pair. Given the pairs, the Fourier response has modulus at most $e^{-X/8}$, where $X$ is their sum of squared chords. A uniform pair has expected squared chord $2D/(d-1)$, so $\mathbb EX=2sD/(d-1)$. Each chord squared is at most four. For two disjoint pairs the product expectation is at most $d(d-1)/[(d-2)(d-3)]$ times the squared single-pair mean, by dropping the disjointness restriction in the numerator. Thus $\mathbb E X^2\le C(\mathbb EX+(\mathbb EX)^2)$. By $1-e^{-x/8}\ge c x/(1+x)$ and Cauchy--Schwarz,

\begin{equation*}
1-|f_s|\ge c\,\mathbb E\frac{X}{1+X}
\ge c\,\frac{(\mathbb EX)^2}{\mathbb E(X+X^2)}
\ge c'\min(sD/d,1).
\end{equation*}

The zero-mean case is immediate. For $sD/d$ above a fixed large constant use the Cauchy estimate and absorb its prefactor. Below that constant use this last lower bound on the loss, with another absolute constant. This proves (G1), and hence $|\prod f_s|\le e^{-c kD}$ globally. It uses no restriction on the degree product beyond the eventual $s_{\max}\le d/2$.

\subsection{Large dispersion}

Let

\begin{equation*}
B_{\vec z}(\tau)=\sum_i|z_i-e^{i\tau}|^2,\qquad
\min_\tau B_{\vec z}(\tau)=2d(1-|\overline z|).
\end{equation*}

Thus $D\ge\tfrac12\min_\tau B_{\vec z}(\tau)$. On $D>L$ reserve half of the bound $e^{-c kD}$ for a factor $e^{-c' kL}$. For the remaining half, approximate the minimizing common phase by its nearest point on a uniform circle grid of $d$ angles. The squared-chord sum increases by at most $C/d$, because its increase is $2d|\overline z|(1-\cos\Delta)$ with $|\Delta|\le\pi/d$. This also holds by choosing any common phase when $\overline z=0$. Since $k=O(d)$, discretization changes the remaining bound by only a bounded factor. At each grid point the resulting product of one-coordinate chord factors integrates to at most $(C/\sqrt k)^d$ by independence. Union over the grid gives

\begin{equation*}
\int_{D>L}\left|\prod f_s\right|\le C d\, e^{-c'kL}(C/\sqrt k)^d.
\end{equation*}

The ratio to $G_d$ is $\exp[-c'kL+O(d)]$, using $\Lambda/k\to1$ and $k/d\to\beta$. Choosing $L$ sufficiently large makes this exponentially small in $d$.

\subsection{Bounded dispersion and an exceptional coordinate}

Fix this $L$ and suppose $D\le L$. For large $d$ the mean is nonzero and close to unit modulus. Write $u_i=z_i/\overline z-1$, so $\sum_i u_i=0$, $\sum_i|u_i|^2=O(L)$, and $\max_i|u_i|\le3$. The exact subset expansion is

\begin{equation*}
f_s=\overline z^{\,s}\sum_{j=0}^s\frac{(s)_j}{(d)_j}e_j(\vec u).
\end{equation*}

To bound its coefficients, take a fixed circle $|v|=1/8$ in $\prod_i(1+v u_i)$. The logarithm has no linear term and satisfies

\begin{equation*}
\left|\sum_i \log(1+v u_i)\right|
\le C |v|^2\sum_i|u_i|^2\le C_L.
\end{equation*}

Cauchy's coefficient bound gives $|e_j|\le C_L8^j$, and $e_1=0$. Since $(s)_j/(d)_j\le(s/d)^j$ and $s_{\max}/d=o(1)$, the sum of the terms with $j\ge2$ is $O_L((s/d)^2)$. Multiplication over rows gives

\begin{equation}
\left|\prod f_s\right|
\le |\overline z|^{dk}\exp[O_L(k s_{\max}/d)]
\le \exp[-kD/2+O_L(k s_{\max}/d)].
\tag{G2}
\end{equation}

The overhead in (G2) is $e^{o(d)}$. The coefficient $1/2$ in its quadratic loss is needed for the exceptional-coordinate argument.

For $D\le L$, the relation to the optimal chord sum sharpens to $D\ge(1-C_L/d)\min_\tau B_{\vec z}(\tau)$. Set $A=(k/2)(1-C_L/d)$ and again discretize the minimizing phase on the $d$-point grid. The increase in $B_{\vec z}$ costs a bounded factor in (G2). If some coordinate is at circular distance $\eta$ or more from the minimizing phase, it remains at distance at least $\eta/2$ from the nearest grid point for sufficiently large $d$.

A single unflagged coordinate has the integral

\begin{equation*}
J(A)=\int_{-\pi}^{\pi}e^{-A|e^{i\theta}-1|^2}\,\frac{d\theta}{2\pi}
=\frac{1+O(A^{-1})}{\sqrt{4\pi A}}.
\end{equation*}

Indeed the chord square differs from $\theta^2$ by $O(\theta^4)$ and is at least $c\theta^2$. Bounding the difference of exponentials by $C A\theta^4e^{-cA\theta^2}$ gives the stated relative error. At a flagged coordinate with $|\theta|\ge\eta/2$, keep half its exponent to gain $e^{-cA\eta^2}$ and integrate the rest. This yields at most $C J(A)e^{-c'A\eta^2}$ for that coordinate. Union over the possible flags and grid points proves

\begin{equation*}
\int_{\{D\le L,\ \mathrm{exception}\}}\left|\prod f_s\right|
\le C d^2 e^{O_L(k s_{\max}/d)} J(A)^d e^{-c' A\eta^2}.
\end{equation*}

Since $4\pi A=2\pi k(1-O_L(1/d))$, the ratio $J(A)^d/G_d$ is at most $\exp(O_L(1)+O(|\log(dk)|)+O(k s_{\max}/d))$. Here one may also use the simpler bound $e^{o(d)}$ for fixed $\eta$. This proves the exceptional-coordinate estimate for fixed cutoff, or for a shrinking cutoff whose squared size dominates $s_{\max}/d+(\log d)/d$.

In the retained region the angles relative to the mean direction are all smaller than $\eta$ and have squared sum $O(L)$ by the chord bound. Subtracting their arithmetic mean gives a zero-sum angle chart with maximum deviation at most $2\eta$ and squared sum still $O(L)$. For small fixed $\eta$ this chart is injective modulo a common rotation, because pairwise angle differences uniquely determine its small zero-sum lift. The local cumulant expansion and its correction constant concern this retained chart and are separate from the absolute estimates established here.


\subsection{A sharper shrinking cutoff}

Retaining the degree-two term in the centered elementary-symmetric expansion improves the overhead. Since $e_2(\vec u)=-\frac12\sum_i u_i^2$,

\begin{equation*}
\mathrm{Re}\,e_2(\vec u)
=\frac12\sum_i |u_i|^2-\sum_i(\mathrm{Re}\,u_i)^2
\le \frac{D}{2|\overline z|^2}.
\end{equation*}

The terms of degree at least three have sum $O_L((s/d)^3)$ by the same Cauchy estimates. The logarithm of the modulus of the centered factor therefore has upper bound
$[(s)_2/(d)_2]D/(2|\overline z|^2)+O_L((s/d)^3)$, with its quadratic logarithm error absorbed there. Use
$s/d-(s)_2/(d)_2=s(d-s)/(d(d-1))$ and $|\overline z|^{-2}=1+O_L(1/d)$. Summing over rows gives

\begin{equation*}
\log\left|\prod f_s\right|
\le-\Lambda D/2 + O_L\bigl(1+k(s_{\max}/d)^2\bigr).
\end{equation*}

Thus in the grid proof one can use $A=(\Lambda/2)(1-O_L(1/d))$. The exceptional-coordinate loss then dominates its overhead and the polynomial unions whenever $\eta_d^2\gg (s_{\max}/d)^2+(\log d)/d$.

\section{Balancing sparse row supports at a linear column quota}\label{sec:sparse-fixed-profile-edgeworth}

Let $d\to\infty$ and let a finite list of integers $s_1,\ldots,s_n\in[0,d]$ have sum $dk$, where $k$ is an integer with $k/d\to\beta\in(0,\infty)$. Assume $\max_i s_i=o(d)$. Choose independent uniform $s_i$-subsets of $[d]$ as the row supports and write $B$ for the probability that all columns have sum $k$. Set $\nu=(d-1)/2$ and

\begin{equation*}
a_s=\frac{s(d-s)}{d(d-1)},\qquad \Lambda=\sum_i a_{s_i}.
\end{equation*}

Then, uniformly along every such sequence of profiles,

\begin{equation*}
B=\sqrt d\,(2\pi\Lambda)^{-\nu}
\exp\!\left(-\frac{1}{12\beta}\right)(1+o(1)).
\end{equation*}

Zero rows have no effect. One may replace $1/\beta$ in the exponential by $d/k$. No distributional assumption on the row-size list is needed. In particular $\Lambda/k\to1$ since $\max s_i=o(d)$. We keep the exact $\Lambda$ in the power of order $d$.

The smaller range $\max_i s_i=o(d^{1/3})$ already follows from Corollary 5.2 of Greenhill, McKay and Wang by Stirling expansion of the binomial normalizer. \href{https://doi.org/10.1016/j.jcta.2005.03.005}{Version of record}

\subsection{Angular coordinates and localization}

For a real angle vector $y$ of zero sum let $Q_j=\sum_i y_i^j$, $M_y=\max_i|y_i|$, and, for a uniform $s$-subset $I$, $L_s=\sum_{i\in I}y_i$. Pinned-coordinate Fourier inversion for the first $d-1$ columns gives the integral of $\prod_i\mathbb E\exp(iL_{s_i})$ on the quotient by common angular shifts, in zero-sum lifts. The centered integrand is periodic on the projected lattice since the total degree is $dk$. Its Euclidean covolume is $(2\pi)^{d-1}/\sqrt d$, from the determinant $1/d$ of the Gram matrix for the first $d-1$ projected standard basis vectors. Thus the Gaussian approximation with precision $\Lambda$ has normalized mass $B_0=\sqrt d(2\pi\Lambda)^{-\nu}$.

The global localization lemma bounds the absolute contribution outside a common arc of arbitrarily small fixed width and bounded squared norm by $o(B_0)$. More explicitly, first choose a sufficiently large fixed dispersion cutoff depending on $\beta$. Outside it, the integral is negligible. Inside it, the integral on configurations with two coordinates separated by more than any fixed small angle is again negligible. In the remaining common arc take the unique short lifts with zero mean. The chord-dispersion cutoff bounds $Q_2$ by a fixed $L$, since chord and angle distances are comparable in this arc. Conversely the local domain contains a cube of small coordinate width intersected with a ball $Q_2\le L_1$ for a fixed $L_1>1/\beta$, by taking the initial cutoff large enough, since dispersion is at most $Q_2$. Let $\mathcal A$ denote the actual retained local domain. Its indicator lies between the indicators of the inner and outer regions just described, both of asymptotically full Gaussian mass. In the Gaussian expectation below, the integrand is multiplied by $\mathbf 1_{\mathcal A}$. Uniform integrability of the modulus on the outer region makes the omitted Gaussian mass negligible. This comparison uses absolute errors, so applies to the complex integrand. Injectivity in the quotient follows because the difference of pairwise lifted angles between any two such representations has absolute value $<2\pi$. \hyperref[sec:sparse-fixed-profile-global-localization]{Proof of global localization}

\subsection{Uniform logarithmic expansion for one row}

For real $y$ with zero sum, $K_s(z)=\log\mathbb E e^{zL_s}$ has a single-valued analytic branch through $K_s(0)=0$ in $|z|<c/M_y$, for an absolute $c>0$. The case $M_y=0$ is immediate. In that disk each $e^{zy_i}$ has positive real part. For variables $u_i$ in the open right half-plane, $e_s(u)$ is nonzero: the roots of $\prod_i(z+u_i)$ and of all its derivatives lie in the open left half-plane by Gauss--Lucas, and its derivative of order $d-s$ at zero is $(d-s)!e_s(u)$. Apply this with $u_i=e^{zy_i}$. With $s=o(d)$ we can assume $s\le d/2$. For $|z|M_y\le c$, elementary symmetric mean comparison on the positive reals gives

\begin{equation*}
\log|\mathbb Ee^{z L_s}|\le s\log\left(d^{-1}\sum_j e^{(\mathrm{Re}\,z)y_j}\right)
\le C a_s |z|^2 Q_2 .
\end{equation*}

The first inequality is the elementary symmetric mean bound $e_s(u)/\binom ds\le (\sum u_i/d)^s$, also obtainable by the sampling-without-replacement comparison. The second uses $\sum y_j=0$ and $s/d\le2a_s$. Borel--Carathéodory applied on two concentric subdisks bounds the absolute logarithm by $C a_s Q_2/M_y^2$ on $|z|\le c' / M_y$. Cauchy's estimates for the resulting power series show that at $z=i$ its remainder after degree four is $O(a_s Q_2 M_y^3)$, if $M_y\le\eta$ for a sufficiently small absolute $\eta$.

For completeness, with $p_j=(s)_j/(d)_j$ (falling factorials), the second, third and fourth cumulants of $L_s$ are respectively

\begin{equation*}
a_s Q_2,\qquad \gamma_s Q_3,\qquad A_sQ_4+3B_sQ_2^2,
\end{equation*}

where
$\gamma_s=a_s(d-2s)/(d-2)$, $A_s=p_1-7p_2+12p_3-6p_4$, and $B_s=(s)_2(d-s)_2/(d)_4-a_s^2$. These identities follow by expanding the moments using the inclusion probabilities for distinct indices and $Q_1=0$, then subtracting three times the squared second moment. For instance the coefficients of $Q_4,Q_2^2$ in the raw fourth moment are $A_s,3(p_2-2p_3+p_4)$, and the latter parenthesis equals $(s)_2(d-s)_2/(d)_4$ by requiring two prescribed elements inside and two outside a random subset. At positive $s\le\max s_i=o(d)$, uniformly,

\begin{equation*}
\gamma_s/a_s=1+o(1),\qquad A_s/a_s=1+o(1),\qquad d B_s/a_s=-1+o(1).
\end{equation*}

Indeed $p_j\le (s/d)^j$ and $a_s\sim s/d$ in this range. For the last identity write $u=s(d-s)$ and observe that $B_s/a_s=(u-d+1)/[(d-2)(d-3)]-u/[d(d-1)]$. The error after multiplying by $d$ and adding one is $O(1/d+s/d)$. Zero-size rows have zero cumulants.

\subsection{Integrating the correction}

Let $Y_i=(Z_i-\overline Z)/\sqrt\Lambda$ with $Z_i$ independent standard normals, giving the real isotropic Gaussian in the zero-sum hyperplane. Write $\Gamma=\sum\gamma_s$, $A=\sum A_s$ and $B'=\sum B_s$, summing over all rows. The logarithmic correction to the Gaussian in the integral on the retained region is

\begin{equation*}
-\frac{i\Gamma}{6}Q_3+\frac A{24}Q_4+\frac{B'}8 Q_2^2
+O(\Lambda Q_2M_y^3).
\end{equation*}

Under this Gaussian, $Q_2\to1/\beta$ in probability, $\Lambda Q_4\to3/\beta$ in probability, and $M_Y=O_{\mathbb P}(\sqrt{\log(d+1)/\Lambda})$. The first two limits follow by the law of large numbers for the centered independent normal coordinates and $\Lambda/d\to\beta$. Hence the quartic part converges to $3/(24\beta)-1/(8\beta)=0$. The remainder also converges to zero in probability. For the cubic, the centered sum $d^{-1/2}\sum_i(Z_i-\overline Z)^3$ differs from $d^{-1/2}\sum_i(Z_i^3-3Z_i)$ by a term tending to zero in probability: expand and use $\overline Z=O_{\mathbb P}(d^{-1/2})$ and $\sum_i(Z_i^2-1)=O_{\mathbb P}(\sqrt d)$. The iid central limit theorem gives limit $N(0,6)$, the variance being the third Hermite polynomial's second Gaussian moment. Since $\Gamma/\Lambda\to1$, $\Gamma Q_3$ converges to $N(0,6/\beta)$.

The exponentials of the corrections on the retained region are uniformly integrable. The cubic is imaginary. For any fixed doubled constants it suffices by Cauchy--Schwarz to bound truncated exponential moments of $\Lambda Q_4$ and $\Lambda L M_Y^3$ separately, since $|B'|Q_2^2=O_\beta(L^2)$. For the first, use independent $X_i=Z_i/\sqrt\Lambda$ and $Y=X-\overline X$. The event $|\overline X|\le\eta$ has probability tending to one independently of $Y$. On that event and $M_Y\le\eta$, $|X_i|\le2\eta$ and $\sum Y_i^4\le16\sum X_i^4$. With $\eta$ sufficiently small each of the resulting one-dimensional truncated expectations is at most $1+C/\Lambda$: use $e^v-1\le ve^v$ and absorb $C\Lambda x^4$ in a quarter of the quadratic exponent. The product is bounded as $d/\Lambda=O(1)$.

For the second term choose $\eta$ still smaller, depending on $L$, so $C L x^3\le x^2/8$ for $x\le\eta$. Each $Y_i$ has variance at most $1/\Lambda$, hence $\mathbb P(M_Y>x)\le2d e^{-\Lambda x^2/2}$. Below $x_0=\sqrt{K\log(d+1)/\Lambda}$ the exponent $C\Lambda L x_0^3=o(1)$. Above $x_0$, integrating the tail bound against $3 C\Lambda L x^2 e^{C\Lambda L x^3}$ gives $o(1)$ for sufficiently large fixed $K$, also accounting for the boundary term at $x_0$. This proves the asserted bounded exponential moments even after doubling constants.

The retained region has Gaussian probability tending to one. Uniform integrability and convergence in distribution of the displayed correction give the limiting Gaussian expectation of $\exp(-iZ/6)$ for $Z\sim N(0,6/\beta)$. This is $\exp(-1/(12\beta))$. The Haar normalization and the negligible global tail prove the formula. As a check, when all nonzero rows have size one, the exact balance probability is $(dk)!/(d^{dk}(k!)^d)$, and Stirling's formula gives the same exponential correction with $\Lambda=k$.

\section{Sparse-profile Edgeworth transfer for bounded rank ratio}\label{sec:edgeworth-bounded-rank-ratio}

Assume the following uniform sparse fixed-profile Edgeworth estimate. If an $n\times d$ binary array is formed by independent uniform subsets of prescribed row sizes $s_j$, $\sum_j s_j=di$, $i/d\to\beta\in(0,\infty)$, and $\max_j s_j=o(d)$, the probability of all column sums equal to $i$ is

\begin{equation*}
\sqrt d(2\pi\Lambda)^{-\nu}\exp[-1/(12\beta)+o(1)],
\quad \Lambda=\sum_j\frac{s_j(d-s_j)}{d(d-1)},\quad\nu=(d-1)/2.
\end{equation*}

The uniform error is needed across all retained profiles; this note is conditional on that input.

Let $d\to\infty$, $b=d^2/n\to\infty$, and $b=o(d)$. At $k/d\to\beta\in(0,\infty)$, work at fixed feasible integer offsets $i=k+O(1)$. Let $m=di/n\asymp b$, $H(z)=\sum_{s\le r}\binom ds z^s$, $A_i=[z^{di}]H(z)^n$, and $S$ be the capped binomial power tilt of mean $m$, with variance $V$. At exact saturation use the deterministic row size $r$. Set $\Lambda_*=i(d-m)/(d-1)$. If $r=O(b)$, all profiles with total $di$ have $s_j\le L=r=O(m)=o(d)$. More generally the transfer below also holds if $L=A(m+\log n)=o(d)$ and both the actual balanced arrays and the scalar law conditioned on its total have $o(1)$ mass above this maximum row size, with the scalar unconditioned second-moment clipping error $o(m)$. The latter case is proved by a switch for $r/b\to\infty$ in \hyperref[sec:edgeworth-finite-lower-overlap]{the lower-quota overlap note}.

\subsection{Averaging the Edgeworth covariance}

With $\widehat V=n^{-1}\sum_j(s_j-m)^2$ on profiles of total $di$,

\begin{equation*}
\Lambda=\Lambda_*\left(1-\frac{\widehat V}{m(d-m)}\right).
\end{equation*}

On profiles with all rows at most $L$, one has $\widehat V\le m(L-m)$, and $L=o(d)$. Their covariance multiplier $(\Lambda_*/\Lambda)^\nu$ is thus at most $\exp(CL)$. The row size law under independent scalar rows conditioned on total $di$ has profile weights proportional to $\prod_j\binom d{s_j}$. Conditioning costs at most $nr+1$: the total size mass function is log-concave, and Jensen for its concave log interpolation at its integer mean bounds the log mass below by minus the entropy, at least $-\log(nr+1)$. For saturation conditioning is trivial.

We need probability errors smaller than $e^{-CL}/(nr+1)$. For independent rows use $Y=(S-m)^2\mathbf1_{\{S\le L\}}/m$, with $Y=(S-m)^2/m$ without clipping when $L=r$. The standard cap variance comparison $V\le m(1-m/d)\le m$ applies after every real tilt. Writing $m(t)$ for the mean after tilt by $e^{tS}$ gives $0\le m'(t)\le m(t)$, hence the centered log mgfs are bounded by $m(e^{\pm t}-1\mp t)$ for $t\ge0$. The upper bound follows from $m(t)\le m e^t$; the lower-tail bound follows from $m(-t)\ge m e^{-t}$. Chernoff gives

\begin{equation*}
\mathbb P(|S-m|\ge x)\le 2\exp[-c x^2/(m+x)].
\end{equation*}

If $m\ge\log n$ or if $L=r=O(m)$, then $L\le C m$. On the nonzero support of $Y$, $|S-m|\le C'm$, so the last inequality gives $\mathbb P(Y>u)\le2e^{-c' u}$ uniformly. Thus $Y$ has a uniformly bounded exponential moment near zero, and the centered log mgf is $O(t^2)$ there by Taylor and that bound. The row average deviates from its expectation by more than $\varepsilon=((L+\log n)/n)^{1/4}$ with probability at most $2e^{-c n\varepsilon^2}$. Here $\varepsilon\to0$ and $n\varepsilon^2/(L+\log n)\to\infty$ since $L=o(d)$ and $n/d\to\infty$. If instead $m\le\log n$ and $L=O(m+\log n)$, use $Y=O(\log^2 n)$ because $m\to\infty$. Hoeffding at tolerance $n^{-1/8}$ yields the same required superexponential comparison with $L+\log n$. In either case conditioning on total and multiplying by $\exp(CL)$ still gives $o(1)$ exceptional contribution. The near-saturation deterministic law also satisfies the same bounds.

On typical retained profiles $\widehat V=V+o(m)$, by the clipping error assumption. There $\widehat V=O(m)$, and Taylor gives $\nu\log(\Lambda_*/\Lambda)=\nu V/[m(d-m)]+o(1)$, since the square term is $O(d^{-1})$. Their conditional probability tends to one. The uniform Edgeworth estimate now gives

\begin{equation*}
a_i=A_i\sqrt d(2\pi\Lambda_*)^{-\nu}
  \exp\!\left[-\frac1{12\beta}+\frac{\nu V}{m(d-m)}+o(1)\right].
\end{equation*}

If clipping was used, the count of actual balanced arrays discarded is a relative $o(1)$ by assumption, so the formula holds for $a_i$ itself. The variance correction is bounded by $1/2$ up to a vanishing error. The auxiliary determinant term has adjacent first difference
$-\nu(\log\Lambda_*(i+1)-\log\Lambda_*(i))=-\nu/i+o(1)$, since $i\asymp d$ and $d/n=o(1)$. Its second difference is therefore $o(1)$.

\subsection{Signs at bounded cap ratio}

Let $D=\lfloor nr/d\rfloor$. The \href{https://github.com/ProgramComputer/open-math-research/blob/78e9ef2c0c0b7a5daed126efbe12a46d030d8200/topics/balanced-arrays/orientation-normalization.md}{orientation inequality} gives $\Gamma_i=a_i^2/(a_{i-1}a_{i+1})\ge(1+1/i)^d(1+1/(n-i))^d$, hence strictly ordered geometric nodes $t_i=\sqrt{a_{i-1}/a_{i+1}}$, uniform positive log curvature when $D/d=O(1)$, and correct sign at each $\Gamma_i>4$. The last implication is the elementary pairwise alternating-tail bound. Possible failing internal nodes must have $k/d$ bounded away from zero by orientation and bounded above by $D/d\le r/b$. Take a subsequence with $k/d\to\beta>0$ and $V(k)/b\to v\in[0,\infty)$.

For real quota $x$ use the tuned tilt with mean $dx/n$ and set $E(x)=n\log H(z)-dx\log z$, with $E''=-b/V$ in the interior. Then $A_i=e^{E(i)}L_i$, with $L_i$ the integer-mean lattice probability. The \href{https://github.com/ProgramComputer/open-math-research/blob/78e9ef2c0c0b7a5daed126efbe12a46d030d8200/topics/balanced-arrays/dense-regime/rank-dependent/small-variance.md}{log-concave lattice estimates} give $|dV/dm|\le C(1+\sqrt V)$ and $L_i=(2\pi nV(i))^{-1/2}(1+o(1))$ whenever $n\min(1,V(i))\to\infty$. The first bound makes $\sqrt{1+V}$ Lipschitz; unit quota shifts change the mean by $b/d=o(1)$. The deficit $J=r-S$ is log-concave relative to the geometric law of its mean $h=r-m$ (its ratio to the geometric masses is log-concave), so $V\le h(1+h)$. This follows by subtracting the line through the two crossings of the mass difference when testing any convex function, using equal masses and means. The \href{https://github.com/ProgramComputer/open-math-research/blob/04564a728c21eab6ce511b98a48ccaac754648ff/topics/balanced-arrays/very-tall/very-tall-interior-saddle.md}{modal-atom estimate} gives $V\ge c\min(1,m,h)$.

\begin{itemize}
\item If $v>0$, then $h\gtrsim\sqrt b$, so the upper support distance $nh/d$ diverges, and the lower distance is $k\asymp d$. At bounded offsets $V/b\to v$ uniformly, the lattice probabilities have relative ratios $1+o(1)$, and the determinant and bounded variance correction have curvature $o(1)$. Consequently $\log\Gamma_{k+j}\to1/v$ at every fixed offset.
\item If $v=0$ and $h\ge\eta>0$, then $V$ stays bounded below at neighboring shifts. The Lipschitz bound for $\sqrt{1+V}$ gives relative variation $o(1)$ even when $b\gg d^{1/3}$: $\sqrt{1+V(x)}$ changes by $O(b/d)=o(1)$. Thus the lattice prefactors have log differences $o(1)$ and $V=o(b)$ uniformly across the unit neighborhood. Entropy curvature diverges, while the fixed-margin corrections are bounded, so $\Gamma_k\to\infty$.
\item If $h\to0$, set $e=nh\ge d$ at an internal quota. The \href{https://github.com/ProgramComputer/open-math-research/blob/04564a728c21eab6ce511b98a48ccaac754648ff/topics/balanced-arrays/very-tall/very-tall-upper-quota.md#the-last-od-quotas}{scalar negative-binomial deficit bound} gives positive log curvature of $A_i$ at least $(1-o(1))b/(h+b/d)\to\infty$, including the case with an exactly saturated neighbor. The fixed-margin formula above is valid at saturation, and its prefactor and bounded correction preserve the divergence. Thus $\Gamma_k\to\infty$.
\end{itemize}

In the first case, fixed-offset magnitudes normalized by the central term at $-t_k$ tend to $e^{-j^2/(2v)}$. The global orientation curvature is bounded below by a positive constant since $D/d=O(1)$, giving a summable Gaussian majorant. The resulting alternating theta sum is strictly positive by Gaussian Poisson summation. The divergent cases directly give correct signs. Thus every internal geometric sign is eventually correct in this regime, and the full root count consists of simple negative zeros. No assumption $b=o(d^{1/3})$ is used.

\section{Finite lower-quota overlap below quadratic height via sparse Edgeworth}\label{sec:edgeworth-finite-lower-overlap}

Assume the uniform sparse fixed-profile Edgeworth input in \hyperref[sec:edgeworth-bounded-rank-ratio]{the covariance-transfer note}. Let $d\to\infty$, $b=d^2/n\to\infty$, $b=o(d)$, $r/b\to\infty$, $1\le r<d$, and $k/d\to\beta\in(0,\infty)$. For $a_i$ the original capped counts, put $\Gamma_i=a_i^2/(a_{i-1}a_{i+1})$, $t_k=\sqrt{a_{k-1}/a_{k+1}}$, and let $w_x$ be the positive occupied tilt with mean $dx/n$. Let $v_k$ be the first positive hitting point of phase $k\pi$ in the \hyperref[sec:analytic-angular-signs]{angular convention}. Then, for every fixed integer $j$,

\begin{equation*}
\log\Gamma_{k+j}\to1/\beta,\qquad
\log t_k=d\log w_k+\frac{d}{2k}+o(1),\qquad
d\log v_k=\log t_k+o(1).
\end{equation*}

In particular both the geometric and angular test points give the correct sign $(-1)^k$. No height assumption stronger than $n/d\to\infty$ is used within these sequences.

\subsection{Switching and scalar clipping}

At bounded offsets $i=k+O(1)$ put $m=di/n\sim\beta b$, $M=di$, and $L=\lceil B(m+\log n)\rceil=o(d)$ for sufficiently large fixed $B$. Both distances to the quota endpoints diverge since $D/d\sim r/b\to\infty$. Under uniform balanced capped arrays let $N_s$ count those with $s$ ones in row one. Move one of its ones downward in the same column to a nonfull row with no one in that column. There are at least $s(n-i-M/r)$ allowed moves because at most $i$ rows are occupied in the column and at most $M/r$ rows are full. Each array with first row size $s-1$ has at most $(d-s+1)i$ possible reverse moves. Thus $N_s\le(2m/s)N_{s-1}$ for large $d$, since $i/n=m/d=o(1)$ and $m/r=o(1)$. The factorial bound $N_s/\sum_jN_j\le(2m)^s/s!$ and a union bound over rows give $\mathbb P(\max s_j>L)=o(1)$, with an arbitrarily large fixed inverse-polynomial bound by increasing $B$.

For the scalar tuned rows, the Poisson Bernstein bound from the covariance-transfer proof, at $L\gg m+\log n$, removes any prescribed polynomial tail even with fixed polynomial weights. Conditioning on the total, which has probability at least $1/(nr+1)$, therefore leaves $o(1)$ tail probability. The clipping error in $\mathbb E[(S-m)^2]$ is $o(m)$, for example by the support bound $d$ and arbitrary polynomial tail. Applying the covariance transfer on these retained profiles yields

\begin{equation*}
a_i=A_i\sqrt d(2\pi\Lambda_*(i))^{-\nu}
  \exp[-1/(12\beta)+\nu V_i/(m(d-m))+o(1)]
\end{equation*}

where $A_i=[z^{di}](\sum_{s\le r}\binom ds z^s)^n$, $\Lambda_*(i)=i(d-m)/(d-1)$, $\nu=(d-1)/2$, and $V_i$ is the scalar size variance. This is where the Edgeworth input replaces Greenhill--McKay--Wang and removes the earlier $b=o(d^{1/3})$ restriction.

At any tilt with $m\asymp b$, the untruncated binomial mean $m_0$ belongs to $[m,2m]$ eventually: the truncated mean at untruncated mean $2m$ is $(2+o(1))m$, because $r/m\to\infty$ and the removed tail has order $o(m^{-A})$ for every fixed $A$. Monotonicity of the truncated mean gives the upper bound. Fixed raw deleted moments satisfy that tail estimate by the binomial exponential moment at a fixed positive argument, and hence $m_0=m+o(1)$, $V_i\sim m$, and central absolute moments through order four are bounded by the usual binomial scales $O(m^{j/2})$ for $2\le j\le4$. These estimates also hold uniformly over bounded real quota shifts. Therefore the variance correction in the displayed coefficient formula tends to $1/2$.

Write $E(x)=n\log H(w_x)-dx\log w_x$. The \href{https://github.com/ProgramComputer/open-math-research/blob/78e9ef2c0c0b7a5daed126efbe12a46d030d8200/topics/balanced-arrays/dense-regime/uniform-lattice-lemma.md}{integer-mean local limit} gives $A_i=e^{E(i)}(2\pi nV_i)^{-1/2}(1+o(1))$. On fixed translates $E'(x)=-d\log w_x$ and $E''(x)=-b/V_x\to-1/\beta$ uniformly. The lattice prefactor has log differences $o(1)$ since neighboring variances have ratio tending to one. The covariance prefactor has log first difference $-\nu/i+o(1)$ and second difference $o(1)$. In the signed centered first difference of $-\log a_i$, the locally constant limiting entropy curvature cancels between the two halves; this gives $\log t_k=d\log w_k+d/(2k)+o(1)$. The positive second log differences tend to $1/\beta$.

\subsection{Phase overlap at every crossing}

For any occupied tilt crossing phase $k\pi$, write $\mu=\mathbb ES$, $y=n\mu/d$, $T=nr/d$, $\varphi=\pi/d$, $G(s)=s(d-s)/[d(d-1)]$, and $a=\mathbb E G(S)$. The global phase bound $|\Psi/\pi-y|\le C d+C\min(y,T-y)/\sqrt d$ from \href{https://github.com/ProgramComputer/open-math-research/blob/4fb624ba937a8eecc1361a6066185cd81ab7b172/topics/balanced-arrays/quadratic/proof.md}{the earlier phase argument} first gives $y=O(d)$, hence $\mu=O(b)$. At this step $\mu$ need not yet be comparable to $b$. The \href{https://github.com/ProgramComputer/open-math-research/blob/78e9ef2c0c0b7a5daed126efbe12a46d030d8200/topics/balanced-arrays/dense-regime/rank-dependent/small-variance.md}{uniform third moment estimates} give $\mathbb E|S-\mu|^3\le C V\sqrt{1+V}$, $V\le\mu$, so the radial phase error multiplied by $n$ is $O(n(1+b)^{3/2}/d^3)=o(1)$ by centering at $\mu$. Even the coarser raw third-moment scale would give $O(b^2/d)=o(d)$, sufficient here. Also $a\asymp\mu/d$ and $\mathbb E[G(S)S]/a\le C(\mu+1)$, using $G(s)\le C s/d$ and $\mathbb E S^2=\mu^2+V$. The characteristic functions of both size laws are within $O((1+b)/d)=o(1)$ of one at $\varphi$ by their first moments. Thus the determinant argument is $O((1+b)/d)$ and costs $O(1+b)=o(d)$ in the phase. It follows that $y=k+o(d)$ and $\mu/b\to\beta$ at every crossing.

Now use the raw binomial-tail comparison above at $\mu\sim\beta b$. Set $Z=S-\mu$ and let $\mathbb E_G$ denote size averaging weighted by $G/a$. Lipschitz continuity of $G$ and the fourth centered moment bound give

\begin{equation*}
\mathbb E_G|Z|^3
\le \frac{G(\mu)\mathbb E|Z|^3+(C/d)\mathbb E|Z|^4}{a}
=O(b^{3/2}).
\end{equation*}

The determinant's linear phase shift is
$\mathbb E_G S-\mu=[(d-2\mu)V-\mathbb E Z^3]/[\mu(d-\mu)-V]\to1$
since $V\sim\mu$, $\mu=o(d)$ and $\mathbb E|Z|^3=O(\mu^{3/2})$. The centered third absolute moments about the means of both laws are $O(b^{3/2})$, so their linear argument expansions have cubic phase errors $O(b^{3/2}/d^3)$. Multiplied by $n$ for the radial phase, this is $O(\sqrt b/d)=o(1)$, and multiplied by $\nu$ for the determinant it is $O(b^{3/2}/d^2)=o(1)$. Thus $\Psi/\pi=y-\tfrac12+o(1)$ and $y(v_k)=k+\tfrac12+o(1)$. Integrating the identity $(d\log w_x)'=b/V_x\to1/\beta$ over this bounded shift proves $d\log v_k=\log t_k+o(1)$.

At every fixed offset $j$ the term magnitude relative to the central term at $-t_k$ tends to $e^{-j^2/(2\beta)}$, with the same limit after the $e^{o(1)}$ radius perturbation. The orientation inequality makes the ratios strictly decrease outward, and the first outward magnitude ratios stay below one. The two alternating tails are bounded by their first omitted terms. Truncate at a fixed distance, take the parameter limit and then let the truncation grow to obtain the positive Gaussian alternating sum by Poisson summation. Hence both nodes have the required sign. Furthermore, this overlap separates adjacent geometric points by $\log(t_{k+1}/t_k)\to1/\beta$, so substitutions of angular points with logarithmic discrepancy $o(1)$ preserve a join here.

\section{Sparse fixed-profile transfer and the complete finite upper band}\label{sec:sublinear-b-generic-upper-transfer}

2026-11-26 \enspace Private research note

Let $d\to\infty$ and $b=d^2/n\to\infty$ with $b=o(d)$. Thus $n/d\to\infty$ and eventually $n\le d^2$. Let $1\le c=d-r$ with $c/b\to\gamma<\infty$, and let $a_k$ be the balanced-array coefficient. Write $T=nr/d$, $D=\lfloor T\rfloor$, $j=D-k$, and $t_k=\sqrt{a_{k-1}/a_{k+1}}$ at an internal quota.

This note assumes the sparse fixed-profile estimate

\begin{equation}
\Pr(\text{equal column sums}\mid\vec s)
=\sqrt d\,(2\pi\Lambda_{\vec s})^{-(d-1)/2}
    \exp\left[-\frac1{12B}+o(1)\right],
\qquad
\Lambda_{\vec s}=\sum_i\frac{s_i(d-s_i)}{d(d-1)},
\tag{U1}
\end{equation}

uniformly over profiles with total $d\ell$, $\ell/d\to B\in(0,\infty)$ and maximum row size $o(d)$. The row subsets in this probability are independent and uniform at their prescribed sizes. Subject to (U1), every internal quota with $j=O(d)$ has the strict sign $(-1)^k$ at $-t_k$. If $j/d\to\beta>0$, the first angular crossing $v_k$ of $k\pi$, and indeed every crossing, satisfies $\log(v_k^d/t_k)=o(1)$ and has the same sign. No degree-product condition enters this transfer.

\subsection{Relative margin summation}

Complement the arrays into holes. Write $\ell=n-k$, $m=d\ell/n=c+h$, $p=\ell/n$, and

\begin{equation*}
G(w)=\sum_{s=c}^d\binom ds w^s,\qquad U_\ell=[w^{d\ell}]G(w)^n.
\end{equation*}

Because $\ell/d=c/b+j/d+O(1/d)$ is bounded, either $\ell/d\to0$ along a subsequence, or we may pass to $\ell/d\to B>0$. In the first case normalized coefficient curvature gives $\log\Gamma_k\ge d\log(1+1/\ell)\to\infty$ for $\Gamma_k=a_k^2/(a_{k-1}a_{k+1})$. This alone gives the geometric sign by the alternating-tail bound below. In the second case $\ell\asymp d$, $m\asymp b$, and $p=o(1)$, including at bounded available shifts.

For $m>c$ let $X$ have the tilted weights $\binom ds w_\ell^s$ on $c\le s\le d$, with $\mathbb E X=m$ and variance $V$. It is an unrestricted binomial tilt with parameter $w_\ell/(1+w_\ell)$ conditioned on size at least $c$. The unrestricted mean is at most $m$, so $dw_\ell=O(m)$, and the mass ratio at the step to size $s$ is at most $Cm/s$. At $m=c$ use the deterministic law at $c$. These bounds give $\mathbb E X^q=O_q(m^q)$, since the ratios are at most one half beyond a large fixed multiple of $m$, and the sizes before then are $O(m)$.

In the actual balanced hole arrays, let $A_s$ count those with specified row size $s$. If $s>c$, move a hole from that row to a different row in the same column. There are exactly $s(n-\ell)$ possibilities, and at most $(d-s+1)\ell$ reverse moves, whether or not the reverse donor can spare its hole. Therefore $s(n-\ell)A_s\le (d-s+1)\ell A_{s-1}$. Its mass ratio is at most $m/[(1-p)s]$. Both the actual and candidate row laws consequently have tails smaller than any prescribed inverse power of $n$ above $L=\lceil A(m+\log n)\rceil$ for sufficiently large fixed $A$. The same holds for fixed moments of their tails. Removing actual arrays with any row above $L$ costs a relative $o(1)$ by the union bound. Note that $L=o(d)$.

Under independent candidate rows, condition on $\mathcal E=\{\sum_iX_i=d\ell\}$. The profile probabilities are proportional to $\prod_i\binom d{s_i}$ with normalizing sum $U_\ell$. The conditioning event has probability at least $1/(nd+1)$: the sum pmf $f$ is log-concave with interval support by convolution, so Jensen at its integer mean on the piecewise-linear interpolation of $\log f$ yields $\log f(d\ell)\ge\mathbb E\log f(\sum_iX_i)\ge-\log(nd+1)$. The deterministic case is immediate. Thus candidate profiles with any row above $L$ also have conditional weight $o(1)$.

For a retained profile put $R=n^{-1}\sum_i(s_i-m)^2$ and $\nu=(d-1)/2$. The exact identity

\begin{equation*}
\Lambda_{\vec s}
=\Lambda_0\left(1-\frac{R}{m(d-m)}\right),\qquad
\Lambda_0=\frac{\ell(d-m)}{d-1},
\end{equation*}

follows by expanding $\sum_i s_i(d-s_i)$. On retained profiles, $R\le Lm$ since $s_i^2\le Ls_i$. Hence the weight $(\Lambda_{\vec s}/\Lambda_0)^{-\nu}$ is at most $e^{O(L)}$, which may be superpolynomial. We next control its conditional expectation.

A binomial law conditioned to be at least $c$ has variance at most its mean, at every real exponential tilt. Indeed its relative likelihood to the binomial law with the same mean is geometric on its interval support. The signed mass difference is positive on an interval and negative outside. Equal total mass and mean allow subtraction of a chord from a convex test function, proving convex-order domination. If $\mu_t$ denotes the mean after real tilt by $e^{tX}$, then $\mu_t'\le\mu_t$. Integration gives $\mu_t\le m e^t$ for $t\ge0$ and $\mu_t\ge m e^t$ for $t\le0$. Integrating these inequalities in their respective directions proves

\begin{equation*}
\log\mathbb E e^{t(X-m)}\le m(e^t-1-t)
\end{equation*}

for every real $t$. Thus $\Pr(|X-m|\ge u)\le2\exp[-c\min(u^2/m,u)]$ by exponential Markov.

Define the clipped variable $W=(X-m)^2/m$ on $X\le L$, and set it to zero otherwise. Its mean differs from $V/m$ by an arbitrarily small inverse power of $n$ by the tail estimate. If $m\ge\log n$ then $L=O(m)$, so the displayed Bernstein-type bound gives $\Pr(W>x)\le2e^{-c_A x}$ for all $x\ge0$. The range above a fixed multiple of $m$ has probability zero, and below it the two alternatives in the Bernstein exponent are each at least a constant times $x$. Consequently $W$ has a uniform positive exponential moment and its centered log moment generating function is $O(t^2)$ in a fixed interval around zero. With tolerance $\epsilon=(L/n)^{1/4}=o(1)$, exponential Markov gives deviation probability $\exp[-\Omega(n\epsilon^2)]$ for its sample mean. Here $n\epsilon^2=\sqrt{nL}\gg L+\log(nd)$ because $L=o(d)$ and $n/d\to\infty$.

If $m\le\log n$, then $L=O(\log n)$ and $W=O(L^2/m)$. At tolerance $n^{-1/10}$, Hoeffding gives deviation probability $\exp[-\Omega(n^{4/5}m^2/L^4)]$, again with exponent dominating $L+\log(nd)$. In both cases, dividing by $\Pr(\mathcal E)$ and multiplying by the maximal weight $e^{O(L)}$ leaves the bad retained profiles negligible. On the other retained profiles, $R/m=V/m+o(1)$ and $V/m\le1$. Therefore

\begin{equation*}
-\nu\log\left(1-\frac{R}{m(d-m)}\right)
=\frac{(d-1)V}{2m(d-m)}+o(1)
\end{equation*}

uniformly there. Applying (U1) on the retained profiles and using the actual relative tail bound proves

\begin{equation}
a_k=\mathcal K(\ell) U_\ell
\exp\left(\frac{(d-1)V}{2m(d-m)}-\frac1{12B}+o(1)\right),
\qquad
\mathcal K(\ell)=\sqrt d\thinspace(2\pi\Lambda_0)^{-\nu}.
\tag{U2}
\end{equation}

This applies uniformly under bounded available shifts. At saturation $R=V=0$ deterministically, and the same formula follows directly. The correction exponent other than $\log\mathcal K+\log U_\ell$ is bounded. Moreover $\log\mathcal K(\ell+1)-\log\mathcal K(\ell)=-d/(2\ell)+o(1)$, since $\Lambda_0=\ell(d-m)/(d-1)$, the second factor changes logarithm by $O(1/n)$ per shift, and $\ell\asymp d$. The second difference is $o(1)$.

\subsection{Scalar shape and all finite upper signs}

In addition to $V\le m$, the excess $X-c$ has variance at most $h(1+h)$ by the same relative log-concavity and equal-mean convex-order comparison with a geometric law. The log-concave lattice moment inequality $\mathbb E|X-m|^3\le CV(1+\sqrt V)$ and differentiation in log tilt give $|dV/dm|\le C(1+\sqrt V)$. In particular $\sqrt{1+V}$ is uniformly Lipschitz as a function of $m$. A bounded quota shift changes $m$ by $O(b/d)=o(1)$, so it changes $\sqrt{1+V}$ by $o(1)$. Thus $V$ is relatively stable whenever bounded below, even when its absolute variation need not tend to zero. We use the uniform log-concave lattice local limit at an integer mean, which gives conditioning probability $(1+o(1))/\sqrt{2\pi nV}$ whenever $V$ is bounded below. \href{https://github.com/ProgramComputer/open-math-research/blob/78e9ef2c0c0b7a5daed126efbe12a46d030d8200/topics/balanced-arrays/fifth-thirds/subquadratic-scaling-corank-upper.md#scalar-limits-and-stability-at-quota-shifts}{Log-concave local limit and centered-moment estimates}

Write $E(\ell)=n\log G(w_\ell)-d\ell\log w_\ell$ for the interior real tilt exponent. The tilt equation gives

\begin{equation*}
E'=-d\log w_\ell,\qquad E''=-b/V.
\end{equation*}

Suppose first that $j/d\to\beta>0$, so $h/b\to\beta$ and $m/b\to\gamma+\beta=B$. Trial an unrestricted binomial mean $c+\beta b/2$. Its lower tail below $c$ is $e^{-\Omega(b)}$ by Chernoff, since its separation from $c$ is a fixed fraction of its order-$b$ mean. Conditioning changes its mean by $O(be^{-\Omega(b)})$, so the resulting mean is less than $m$. Monotonicity gives the same order-$b$ separation for the actual unrestricted mean. Conditioning at that mean changes raw moments of each fixed order $q$ by $O_q(b^qe^{-\Omega(b)})$. Hence $V\sim m$ as the unrestricted binomial parameter tends to zero, uniformly at bounded real shifts. The local limit expresses $U_\ell=e^{E(\ell)}(1+o(1))/\sqrt{2\pi nV}$. Both its prefactor and the correction in (U2) vary by $o(1)$ in log at fixed shifts. The centered first difference of $E$ is $E'+o(1)$, since $E''$ has a common finite limit on both halves. Together these give

\begin{equation}
\log\Gamma_{k+i}\longrightarrow 1/B\quad(i\text{ fixed}),\qquad
\log t_k=-d\log w_\ell-d/(2\ell)+o(1).
\tag{U3}
\end{equation}

If $j/d\to0$ and $\ell/d\to B>0$, then $B=\gamma$. Pass to subsequences in the bounded quantities $V/b$ and, where needed, $h$.

If $V/b\to v_*>0$, then $V\le h(1+h)$ forces $h\gtrsim\sqrt b$ and $j\gtrsim d/\sqrt b\to\infty$, so all fixed neighboring quotas exist. Relative variance stability, (U2), and the local limit give $\log\Gamma_{k+i}\to1/v_*$ at each fixed offset. Here the correction in (U2) has a common finite limit at each shift.

If $V/b\to0$ and $h\ge\eta>0$, then $V$ is uniformly bounded below. Otherwise Chebyshev at very small variance would force one atom to have mass above $0.9$. At an interior mode $s$, the product of its two adjacent mass ratios is $s(d-s)/[(s+1)(d-s+1)]\ge1/4$, which precludes this. At the upper endpoint $d$, such a mass contradicts $m=O(b)=o(d)$. At the lower endpoint $c$, if the first forward ratio is $q$, log-concavity and the large modal mass give $h\le Cq$ and $V\ge c_1q$, contradicting $h\ge\eta$ for sufficiently small variance. This lower bound holds also on bounded real shifts, since they change $h$ by $o(1)$. Relative stability and the local limit give log-prefactor differences $o(1)$, while $E''=-b/V\to-\infty$ uniformly on the shifts. The correction in (U2) again varies by $o(1)$, so $\Gamma_k\to\infty$.

If $h\to0$, put $e=nh\ge d$ at an internal quota. Apply the scalar negative-binomial curvature bound to the excess coefficients $\binom d{c+s}$. It yields

\begin{equation*}
\log\frac{U_\ell^2}{U_{\ell-1}U_{\ell+1}}
\ge \frac{d(n-1)}{n+e+d-1}\log(1+d/e)
\ge (1-o(1))\frac{b}{h+b/d}\longrightarrow\infty.
\end{equation*}

The variances at the three quotas tend to zero by $V\le h(1+h)$, including a possible saturated neighbor. The small changes of the prefactor and correction in (U2) preserve the divergence. \href{https://github.com/ProgramComputer/open-math-research/blob/04564a728c21eab6ce511b98a48ccaac754648ff/topics/balanced-arrays/very-tall/very-tall-upper-quota.md#the-last-od-quotas}{Scalar bound}

For an actual coefficient vector, ratios decrease strictly by normalized log-concavity. At $t_k$ both first outward magnitudes relative to the central term are $\Gamma_k^{-1/2}$ and the magnitudes decrease outwards. Divergent $\Gamma_k$ therefore gives the sign $(-1)^k$ by the first-term bound on each alternating tail. If $\log\Gamma_{k+i}\to C>0$ for all fixed offsets, the normalized magnitude at offset $i$ tends to $e^{-Ci^2/2}$. Cut the alternating tails with error at most the first omitted term. Passing first to the quota limit and then to an increasing truncation gives the positive limit $\sum_{i\in\mathbb Z}(-1)^i e^{-Ci^2/2}$ by the Gaussian Fourier series. The same proof applies at $t_ke^{o(1)}$, since fixed offsets have the same limits and both first outward ratios stay strictly below one. This proves the strict signs along every sequence in the finite upper band.

\subsection{Matching the angular points at positive offset}

For an occupied tilt $v$ let $X$ be the hole size under $w=1/v$, $S=d-X$, $\mu=\mathbb ES$, $y=n\mu/d$, $\phi=\pi/d$, and $\lambda(X)=X(d-X)/[d(d-1)]$. The angular phase is

\begin{equation*}
\Psi(v)=n\arg H_{\rm occ}(ve^{i\phi})-\nu\arg a_\phi,\qquad
a_\phi=\frac{\mathbb E[\lambda(X)e^{-i\phi X}]}{\mathbb E e^{-i\phi X}} ,
\end{equation*}

where $H_{\rm occ}(v)=\sum_{s=0}^r\binom ds v^s$. The occupied sums have their arguments in the sector $[0,r\phi]$, so the phase is continuous on positive tilts. Its limits at zero and infinity are $-\nu\phi$ and $T\pi$, giving ordered first crossings.

At any crossing $\Psi=k\pi$, the general bound $|\Psi/\pi-y|\le Cd+C(T-y)/\sqrt d$ gives $T-y=O(d)$, since $T-k=O(d)$. For this bound use $V\le\min(d/4,h(1+h))\le C\sqrt d\,h$ from convex-order comparisons with the mean-matched binomial and geometric, and use the cubic sine remainder on the centered size characteristic function to bound the row-angle error by $O(V/d^2)$. The curvature argument has absolute value at most $\pi$ from its two sectors. Thus at the crossing the hole mean $m=d-\mu=c+d(T-y)/n$ is $O(b)$, although we have not yet separated the hole mean from the cap.

Set $Z=X-m$. For the hole tilt at these crossings, the centered log-concave moment bound gives $\mathbb E|Z|^3\le CV(1+\sqrt V)\le Cm^{3/2}$. The centered row-angle remainder multiplied by $n$ is $O(nm^{3/2}/d^3)=O(\sqrt b/d)=o(1)$. For the curvature argument, center the phase in both numerator and denominator at $m$. We have

\begin{equation*}
\mathbb E\lambda(X)=m/d+O(m^2/d^2),\quad
\frac{\mathrm{Cov}(\lambda(X),X)}{\mathbb E\lambda(X)}
=\frac Vm+O(m/d).
\end{equation*}

The second identity uses $\mathrm{Cov}(X^2,X)=2mV+\mathbb E Z^3$ and $V\le m$. Crucially the relative weighted second centered moment is $O(m)$, since

\begin{equation*}
\mathbb E[\lambda(X)Z^2]\le \frac{C}{d}\left(mV+\mathbb E|Z|^3\right)
=O(m^2/d).
\end{equation*}

Thus Taylor of the centered weighted characteristic quotient gives $\arg a_\phi=-(V/m)\phi+O(\phi m/d+\phi^2m)=-(V/m)\phi+o(\phi)$. This uses only $m=O(b)=o(d)$, with $m\ge c\ge1$. Consequently every crossing satisfies

\begin{equation*}
\Psi/\pi=y+\frac{V}{2m}+o(1),\qquad y=k+O(1).
\end{equation*}

If $j/d\to\beta>0$, then at every crossing $h/b=(T-y)/d\to\beta$. The trial-binomial calculation now yields $V/m\to1$ there, so $y=k-1/2+o(1)$. At $v_{\rm geo}=t_k^{1/d}$ the drift formula in (U3) shifts the occupied log tilt from $-\log w_\ell$ by $-1/(2\ell)+o(1/d)$. Since $dy/d(d\log v)=V/b\to B$ uniformly on this short interval, this gives the same value $y(v_{\rm geo})=k-1/2+o(1)$. The uniformity follows from $V\le m$ controlling the change of $m$ over an $O(1/d)$ log-tilt interval. Strict increase of $y$ traps every value between this point and any crossing in that same $o(1)$ neighborhood. Consequently $h/b\to\beta$ and $V/b\to B>0$ throughout the connecting interval, so the mean value theorem gives $\log(v_k^d/t_k)=o(1)$. This proves the first crossing statement as well as the perturbative strict sign there.

At an overlap with $j/d\to\beta>0$, the adjacent geometric log spacings tend to $1/B$ by (U3), since the spacing is $(\log\Gamma_k+\log\Gamma_{k+1})/2$. Replacing the last inner test point by its first angular crossing preserves the order. This is the finite upper-band transfer conditional on (U1); the sparse fixed-profile constant and all remaining bands are separate inputs.

\end{document}
